% a06_BALtorus_d3b: continuation of a06_BALtorus_d3 (input from sec/a06_BALtorus_c.tex): the step k+1 as Ba2 treats it, (DOM)_k, and the bad localisations of Ba3 p. 422.

\begin{lemma}[{the new fluctuation at step $k+1$, as \cite[pp.~575--581]{Ba2} treats it}]\label{lem:a06d:K3}
At step $k+1$, \cite{Ba2} writes $\theta_{k+1}A^{(k)}=\theta_{k+1}A'^{(k)}+\tilde B^{(k)}+\theta_{k+1}B^{(k+1)}$ \cite[(2.88)--(2.89)]{Ba2}, $\theta_{k+1}$ a smooth cutoff of $B^k(\Lambda^{(k)}_2)$ changing in a slice of thickness $M$, keeps $(1-\theta_{k+1})A^{(k)}$ in the background $\tilde B$ \cite[p.~579]{Ba2}; expands the $Z^{(j)}$ factors \cite[(2.92)]{Ba2}, $j<k$, on their domains $B^j(\Lambda^{(j)}_2)$ with restrictions $\Lambda^{(j)}_5$, moves them to $B^k(\Lambda^{(k)}_0)$, composes them by \cite[(2.105)]{Ba2} into the polynomial $Q^{(k)}$ \cite[(2.106)]{Ba2}; drops the terms with vertices localised in $B^k(\Lambda^{(k)c}_7)$ and relocalises with a smooth $g^\sharp$ ($\supp g^\sharp\subset\Lambda^{(k)}_7$, $g^\sharp=1$ on $\Lambda^{(k)}_8$) \cite[p.~579]{Ba2}; replaces the lines by $G_k(B^k(\Lambda^{(k)}_2),B^{(k+1)})$ \cite[Prop.~2.6, (2.107)]{Ba2}; and integrates $A',\phi'$ conditionally on $\Lambda^{(k)c}_5$ with $C^{(k)}_{\Lambda^{(k)}_5}$ \cite[(2.115)]{Ba2}. At zero $(k+1)$-fields (so $\tilde B^{(k)}=0=\theta_{k+1}B^{(k+1)}$) write $\mathrm{out}^\theta_k$ for the vacuum generated before the drop, and $\mathrm{out}^\theta_k=\mathrm{out}^\sharp_k+D_k$, where $D_k$ collects the terms with at least one new leg carrying $(1-g^\sharp)\theta_{k+1}$ and $\mathrm{out}^\sharp_k$ is kept in the action; the legs of the new-kernel term $\mathcal K_k$ of \cref{lem:a06d:V6} count as new legs. Then, with $\mathrm{out}_k$ of \cref{lem:a06d:V6} on $\mathfrak D'=(\Omega^{(k)},W,g',c')$, $W\defeq B^k(\Lambda^{(k)}_7)$,
\[
\mathrm{out}^\sharp_k=\mathrm{out}_k[\mathfrak D';C_{\mathrm{act}}]+X_0+X_3,
\]
with $|X_0|\le Ct_ke^{-\delta(5r(t_k)-M-c_d)}N_k$ (the domain moves and the conditioned background $\tilde B=\sum_{l\le k}(1-\theta_{l+1})\theta_lA^{(l)}$ at zero new fields, every slice of which lies in $B^{l-1}(\Lambda^{(l-1)}_2)\setminus B^l(\Lambda^{(l)}_2)\subset B^k(\Lambda^{(k)}_2)^c$, at distance $\ge5r(t_k)-M-c_d$ from $W$), $|X_3|\le Ct_k(2D+5)^3|\mathrm{Col}_k|_k$ with $D\defeq\max(M,D^\sharp)$, $D^\sharp$ the transition width of $g^\sharp$ (families with one new leg carrying $(g^\sharp-g')$, supported within $D$ of $\partial W$), and $|D_k|\le Ct_k(2+(2D+3)^3)|\mathrm{Col}_k|_k$ (the support of $(1-g^\sharp)\theta_{k+1}$ lies in $B^k(\Lambda^{(k)}_2\setminus\Lambda^{(k)}_7)\cup(\theta\text{-slice})\cup\{g^\sharp<1\}$, the first two in $B^k(\mathrm{Col}_k)$ since $\Lambda^{(k)}_2\subset\Lambda^{(k)}_0\subset\Lambda^{(k-1)'}_7$ and $\dist(\Lambda^{(k)}_2,\Lambda^{(k-1)'c}_7)>2r(t_k)>M$). In all three, $t_k(2D+5)^3\le C't_k^{\kappa_0}$, $\kappa_0=1/5$, even for $D\sim r(t_k)$.
\end{lemma}

\begin{proof}
All three are differences of grouped families (\cref{lem:a06d:K2}) with one designated object: a deep domain or background change ($X_0$: the Combes--Thomas bound for the background resolvent difference, (I2a) for domains, (I2b) for the $\Lambda^{(j)}_5$ restrictions; $W$ is at depth $\ge5r(t_k)-c_d$ in $B^k(\Lambda^{(k)}_2)$ and nested inside every $B^j(\Lambda^{(j)}_2)$, $\Lambda^{(j)}_5$), or a smooth multiplier $(g^\sharp-g')$ or $(1-g^\sharp)\theta_{k+1}$ on one new leg ($C^2$-bounded, (I1d)). Bound each by \cref{lem:a06d:V9} and sum with \cref{lem:a06d:V11,lem:a06d:V8}(f); $\theta_k=1=\theta_{k+1}$ on $\supp g^\sharp$, so $g^\sharp\theta_k\theta_{k+1}=g^\sharp$. The last claim is $t(1+\log t^{-1})^{3\varrho}\le Ct^{1/5}$.
\end{proof}

\begin{theorem}[$(\mathrm{DOM})_k$, $0\le k<K$]\label{thm:a06d:V14}
Ba{\l}aban's step $k+1$ (\cite[p.~412]{Ba3}, \cite[pp.~569--582]{Ba2}) consists of: (O1) the collar removal \cite[(2.57)]{Ba2}, down to a vertex datum $(B^k(\Lambda^{(k)}_7),g',c')$; (O2) the translation and the expansion in the new fluctuation; (O3) the integration of the step-$k$ fluctuations with the covariances $C^{(k)}_{\Lambda^{(k)}_5}(B^k(\Lambda^{(k)}_2),0)$ (scalar) and $C^{(k)}_{\Lambda^{(k)}_5}$ (vector) of \cite[p.~419]{Ba3}, and the cumulant expansion; (O4) the counterterm move $k\to k+1$ of \cite[p.~422]{Ba3}, $E_{K,k}=E_{K,k+1}+\mathcal V^{\mathrm{glob}}_k$ and $\dm_{K,k}\to\dm_{K,k+1}$ with the index-$k$ slice moved into $\dm_{k+1}$, followed by localisation; (O5) the rescaling to scale $k+1$ and the new datum $\mathfrak D_{k+1}$. Let $\mathrm{Vac}^{\mathrm{act}}_{k+1}$ be the field-independent part of degree $1..\bar n$ of the resulting polynomial, in the logarithmic convention with the front factor $e^{-\Efin-E_{K,k+1}}$, the removed collar constant $\mathcal R_k[\mathfrak D_k]-\mathcal R_k[\mathfrak D']$ (\cref{cor:a06d:V13}) and the term $D_k$ of \cref{lem:a06d:K3} being dropped and booked by their absolute values. Then
\[
\mathrm{Vac}^{\mathrm{act}}_{k+1}=\mathcal R_{k+1}[\mathfrak D_{k+1}]+\mathcal E'_k+\mathcal O_k,
\]
where $\mathcal E'_k$, $\mathcal O_k$ are finite sums of localised vacuum families with $|\mathcal E'_k|\le Ct_k(2D+5)^3|\mathrm{Col}_k|_k+Ct_ke^{-\delta c_*r(t_k)}N_k\le C't_k^{\kappa_0}|\mathrm{Col}_k|_k+C't_k^{16/5}N_k$ ($D$ as in \cref{lem:a06d:K3}, $\kappa_0=1/5$), $|\mathcal O_k|\le Ct_k|\Lambda^{(k)c}_7|_k$, every family of $\mathcal O_k$ having a designated vertex in $\Lambda^{(k)c}_7$; the constants do not depend on $\eps$, $k$, the torus or the regions. Equivalently, the field-independent output of step $k$, $\mathrm{out}^{\mathrm{act}}_k-\mathcal V^{\mathrm{glob}}_k$, equals $\mathcal R_{k+1}[\mathfrak D_{k+1}]-\mathcal R_k[\mathfrak D']+\mathcal E'_k+\mathcal O_k$.
\end{theorem}

\begin{proof}
Put $W\defeq B^k(\Lambda^{(k)}_7)$, $\Omega^\natural\defeq B^k(\Lambda^{(k)}_2)$, $\mathfrak D'\defeq(\Omega^{(k)},W,g',c')$, $\mathfrak D^\natural\defeq(\Omega^\natural,W,g',c')$; $\Omega^\natural\subset\Omega^{(k)}$ by the nesting $\Lambda^{(k)}_2\subset\Lambda^{(k)}_0\subset\Lambda^{(k-1)'}_7\subset\Lambda^{(k-1)'}_2$ \cite[p.~570]{Ba2}.

1. After (O1) the vacuum part is $\mathcal R_k[\mathfrak D']$, the front factor unchanged; the dropped constant is bounded by \cref{cor:a06d:V13}.

2. (Input domain.) By \cref{lem:a06d:V7}(a), $\mathcal R_k[\mathfrak D']=\mathcal R_k[\mathfrak D^\natural]+X_1$ and $F_k[\mathfrak D']=F_k[\mathfrak D^\natural]+(\delta\text{-families})$, the designated object $\delta G_k(\Omega^{(k)},\Omega^\natural)$ at depth $\ge5r(t_k)-c_d$ (\cref{lem:a06d:V8}(b)).

3. (Step covariances.) $\mathrm{out}^{\mathrm{act}}_k\defeq\mathrm{out}^\sharp_k=\mathrm{out}_k[\mathfrak D';C_{\mathrm{act}}]+X_0+X_3$ (\cref{lem:a06d:K3}), and $\mathrm{out}_k[\mathfrak D';C_{\mathrm{act}}]=\mathrm{out}_k[\mathfrak D^\natural;C_{\mathrm{comp}}]+X_2$ by \cref{lem:a06d:V7}(a) for the input $F_k$ (step~2, then integrated) and \cref{lem:a06d:V7}(c) for $C_{\mathrm{act}}-C_{\mathrm{comp}}$: scalar $C^{(k)}_{\Lambda^{(k)}_5}(\Omega^\natural,0)-C^{(k)}(\Omega^\natural,0)$, vector $C^{(k)}_{\Lambda^{(k)}_5}-C^{(k)}$ (I2b), every designated object at depth $\ge2r(t_k)-c_d$ from $\Lambda^{(k)c}_5$. By \cref{lem:a06d:V9}(b), (c) and \cref{lem:a06d:V11}, $|X_1|+|X_2|\le Ct_ke^{-\delta(2r(t_k)-c_d)}N_k$.

4. (Exact staging.) \cref{lem:a06d:V6}(c) on $\mathfrak D^\natural$: $\mathcal R_k[\mathfrak D^\natural]+\mathrm{out}_k[\mathfrak D^\natural;C_{\mathrm{comp}}]-\mathcal V_k[\mathfrak D^\natural]=\mathcal R_{k+1}[\mathfrak D^\natural\!\uparrow]$.

5. (The counterterm move.) The front factor splits, $e^{-E_{K,k}}=e^{-E_{K,k+1}}e^{-\mathcal V^{\mathrm{glob}}_k}$ with $\mathcal V^{\mathrm{glob}}_k=\mathcal V_k[\mathfrak D^T]$ (largest line index exactly $k$, full coefficient $\dm$). $\mathrm{Ct}$ does not depend on the propagator domain, so $\mathcal V_k[\mathfrak D^\natural]=\mathcal V_k[\mathfrak D^T_{W,g'}]$. \cref{lem:a06d:V6}(c) on the two torus data (where $\Delta=0$) gives $\mathcal V_k[\mathfrak D^T]=\mathrm{out}_k[\mathfrak D^T]-\mathfrak B_{k+1}[\mathfrak D^T\!\uparrow]+\mathfrak B_k[\mathfrak D^T]$ and the same for $\mathfrak D^T_{W,g'}$, hence
\[
Y_k\defeq\mathcal V^{\mathrm{glob}}_k-\mathcal V_k[\mathfrak D^\natural]=\big[\mathrm{out}_k[\mathfrak D^T]-\mathrm{out}_k[\mathfrak D^T_{W,g'}]\big]-\big[\mathfrak B_{k+1}[\mathfrak D^T\!\uparrow]-\mathfrak B_{k+1}[\mathfrak D^T_{W,g'}\!\uparrow]\big]+\big[\mathfrak B_k[\mathfrak D^T]-\mathfrak B_k[\mathfrak D^T_{W,g'}]\big].
\]
By \cref{lem:a06d:V7}(b) each bracket is a sum of families with a designated vertex in $Z^T\defeq W^c\cup\{g'<1\}$ or in the zone where the coefficient changes; so $Y_k$ needs no bound on index-$k$ slices alone, only Proposition~1 for out-type and contour families. Sort by the cube of the designated vertex: families with that cube in $\Lambda^{(k)c}_7$ form $-\mathcal O_k$, families with it in $W$ (then within $M+1$ of $\partial W$) go into $\mathcal E_k$. By \cref{lem:a06d:V9}(a), (c), (d), \cref{lem:a06d:V11} and \cref{lem:a06d:V8}(f): $|\mathcal O_k|\le Ct_k|\Lambda^{(k)c}_7|_k$ and $|Y_k+\mathcal O_k|\le Ct_k(1+(2M+5)^3)|\mathrm{Col}_k|_k$ ($\mathfrak B_{k+1}$ lives at scale $k+1$, which changes the constants by at most $L^3$).

6. (The new datum.) Compare $\mathfrak D_{k+1}$ with $\mathfrak D^\natural\!\uparrow$. The propagator domain changes from $\Omega^\natural$ to $B^{k+1}(\Lambda^{(k)'}_2)$, a deep change (\cref{lem:a06d:V7}(a), depth $\ge5r-c_d$). The vertex datum changes from $(W,g',c')$ to $(B^{k+1}(\Lambda^{(k)'}_7),g_{k+1},c_{k+1})$, including the coefficient move of \cite[(1.28)--(1.29)]{Ba3}, $c'-c_{k+1}=[\dm_{[k]}(1)-\dm_{[k]}(g_{k+1})]+[\dm_k(g')-\dm_k(g_{k+1})]$ with $\dm_{[k]}$ the index-$k$ slice (kinds (b3) and (b1)), all supported within $D_6\defeq M+LM+2$ of $\partial W$ up to exponential tails. By \cref{lem:a06d:V7}(b), \cref{lem:a06d:V9}(d) and \cref{lem:a06d:V8}(f), $|\mathcal R_{k+1}[\mathfrak D^\natural\!\uparrow]-\mathcal R_{k+1}[\mathfrak D_{k+1}]|\le Ct_k(1+(2D_6+3)^3)|\mathrm{Col}_k|_k+Ct_ke^{-\delta(5r-c_d)}N_k$. (If the new legs carry $g^\sharp$, start from the mixed datum, $g'$ on the old legs and $g^\sharp$ on the new; the difference is $X_3$.) The rescaling is exact.

7. (Assembly.) $\mathrm{Vac}^{\mathrm{act}}_{k+1}=\mathcal R_k[\mathfrak D']+\mathrm{out}^{\mathrm{act}}_k-\mathcal V^{\mathrm{glob}}_k=\mathcal R_k[\mathfrak D^\natural]+X_1+\mathrm{out}_k[\mathfrak D^\natural;C_{\mathrm{comp}}]+X_2+X_0+X_3-\mathcal V_k[\mathfrak D^\natural]-Y_k=\mathcal R_{k+1}[\mathfrak D^\natural\!\uparrow]+X_1+X_2+X_0+X_3-Y_k=\mathcal R_{k+1}[\mathfrak D_{k+1}]+\mathcal E'_k+\mathcal O_k$, with $\mathcal E_k\defeq X_1+X_2-(Y_k+\mathcal O_k)+(\mathcal R_{k+1}[\mathfrak D^\natural\!\uparrow]-\mathcal R_{k+1}[\mathfrak D_{k+1}])$ and $\mathcal E'_k\defeq\mathcal E_k+X_0+X_3$.

8. (The exponents.) $t_k\le t_k^{\kappa_0}$, and $e^{-\delta c_*R(1+\log t^{-1})^\varrho}\le C_\kappa t^\kappa$ for every $\kappa$ gives $te^{-\delta c_*r(t)}\le Ct^{16/5}$, likewise for $X_0$'s factor. ($C_\kappa$ is free of $\eps_0$ but can be very large; it is absorbed in the final constants $E_\pm$.)
\end{proof}

No term of $\mathcal R_k$ is bounded in absolute value except its collar part (\cref{cor:a06d:V13}) and differences with a designated object deep inside or in a zone; the restrictions $\chi$, the non-polynomial remainders and the degree-$0$ normalisations are Step~3(C). The field-dependent analogue (\cite[p.~412]{Ba3}: \cite[(1.5)]{Ba3} after $k+1$ steps on the new domains) follows from \cref{lem:a06d:V6,lem:a06d:V7} in the same way, with Proposition~1 under the weaker restrictions on new scalar fields \cite[p.~421]{Ba3}; its interaction-term half at a nonzero background is \BALint\ (\cref{hyp:k01:BALint}).

\paragraph{The bad localisations of \cite[p.~422]{Ba3}.} ``Passing from $k$ to $k+1$ we have to change \dots\ the terms in $\dm_{K,k}(x)(L^k\eps)^2$ and $E_{K,k}$ for which the operators in the decomposition (1.27) have the indices $j\le k$ and at least one of them is equal to $k$ \dots\ then localize them properly, which means that we have to estimate the terms with bad localizations''; the moved line of index $k$ ``is a regular function of both variables and we can treat it as an external field''. Write $S_{(k)}$ for the slice of index $k$ of \cite[(1.27)]{Ba3}, $G_{k+1}=G_k+S_{(k)}$.

\begin{lemma}[the transfer identity]\label{lem:a06d:D21}
On the torus at zero background, $\dm_{K,k}-\dm_{K,k+1}=\dm_{k+1}-\dm_k=:m_k$, and $m_k$ is the sum of the graphs of \cite[(1.23)]{Ba3} whose lines are all $G_k$ or $S_{(k)}$ with at least one $S_{(k)}$ line; the same holds for $E_{K,k}-E_{K,k+1}$.
\end{lemma}

\begin{proof}
By \cite[(1.27)]{Ba3}, $G_{K,k}=S_{(k)}+G_{K,k+1}$ and $G_{k+1}=G_k+S_{(k)}$, for scalar and vector lines alike. Expand each graph of $\dm_K$ multilinearly with every line written $G_k+S_{(k)}+G_{K,k+1}$: $\dm_{K,k}$ collects the terms with at least one line in $\{S_{(k)},G_{K,k+1}\}$, $\dm_{K,k+1}$ those with at least one line in $G_{K,k+1}$, and their difference is the terms with all lines in $\{G_k,S_{(k)}\}$ and at least one $S_{(k)}$, i.e.\ $\dm_{k+1}-\dm_k$. The counterterm vertices inside these graphs are vertices, not lines (\cref{conv:a06d:LOC}(a)), and are not expanded.
\end{proof}

\begin{lemma}[the bad-localisation terms of the moved mass counterterm]\label{lem:a06d:D22}
At step $k+1$ let the moved terms $m_k$, rescaled to the $\eta$-lattice, be localised by the rules of \cref{conv:a06d:LOC} with $\Omega^+_1\defeq B^k(\Lambda^{(k)}_7)$ and $g^+$ ($\supp g^+\subset\Omega^+_1$, $g^+=1$ at distance $\ge LM$ from $\partial\Omega^+_1$ in scale-$k$ units, \cite[p.~412]{Ba3}); assume $R\ge5M+6$, so that $r-1>LM+4$. Put $\mathrm{Bad}\defeq(\Omega^+_1)^c\cup\{z\in\Omega^+_1:\dist(z,\partial\Omega^+_1)<LM+2\}$ and $b_k(x)\defeq m^{\mathrm{glob}}_k(x)-m^{\mathrm{loc}}_k(x)$ for $x\in\Omega^+_1$, the coefficient change of the \cite[(1.7)]{Ba3} vertex caused by the move. Then: \textup{(a)} any localised family of the step-$(k+1)$ expansion containing one \cite[(1.7)]{Ba3} vertex with coefficient $b_k(x)t^2$ at $x$ is a finite sum of Proposition~1 families, each with one $S_{(k)}$ line treated as an external field and at least one vertex localised in $\mathrm{Bad}$, and $|E|\le O(1)e_k^{d_v}\lambda_k^{d_s}\prod(\text{external norms})\,C_Ke^{-(\delta'/2)(\dist(\square(x),\mathrm{Bad})-2)}e^{-(\delta'/2)d(\{\square\})}$, $\delta'\defeq\min(\delta_0,\delta_K)$, $C_K,\delta_K$ the constants of the slice kernel; \textup{(b)} the standalone vertex with both legs external satisfies $\sum_{x\in\Omega^+_1}\eta^3|b_k(x)t^2||\Phi(x)|^2\le Ct^{1/2}p^2\lambda^{-1/2}[|\Lambda^{(k)}_7\setminus\Lambda^{(k)}_8|+e^{-(\delta'/2)(r-LM-6)}|\Lambda^{(k)}_7|]$, and every family with $W\ge2$ containing the $b_k$ vertex obeys the same bound with $t^{1/2-\beta}$; \textup{(c)} in Ba{\l}aban's form: for $x\in B^k(\Lambda^{(k)}_7\setminus\Lambda^{(k)}_8)$ the size is a power of the rescaled couplings ($O(e_k^2+\lambda_k)$ for the coefficient), for $x\in B^k(\Lambda^{(k)}_8)$ it carries in addition $e^{-(\delta'/2)(r-LM-6)}$, the printed $e^{-\delta_1r}$ with $\delta_1=\delta'/4$ once $r\ge2(LM+6)$.
\end{lemma}

\begin{proof}
$m^{\mathrm{loc}}_k$ is the same sum of graphs as $m^{\mathrm{glob}}_k$ with weight $\mathbf 1_{\Omega^+_1}$ on the \cite[(1.6)--(1.7)]{Ba3} vertices and $g^+$ on the vector legs; telescoping the weights gives terms in which at least one internal vertex $z$ carries $1-\mathbf 1_{\Omega^+_1}(z)$ or $1-g^+\ne0$, so $z\in\mathrm{Bad}$. Each graph of $m_k$ has an $S_{(k)}$ line (\cref{lem:a06d:D21}); cut one in a fixed order and regroup into the renormalised families of the cut graph (as in \cref{thm:a06d:S23}, the slice line playing the role of a cut line; the $S_{(k)}$ line belongs to no divergent subgraph of the cut graph); the slice kernel is regular in both variables with $|S_{(k)}|\le C_Ke^{-\delta_K\dist}$ in the two-variable norm, the printed ``regular function of both variables''. Apply Proposition~1 to the resulting families, joining the components of a disconnected cut graph through the slice decay (\cref{lem:a06d:V10}). The tree contains $\square(x)$ and a cube in $\mathrm{Bad}$, so $d(\{\square\})\ge\dist(\square(x),\mathrm{Bad})-2$, and half of the exponent is kept for the localisation sum. The order is $\ge2$, giving $O(e_k^2+\lambda_k)=O(t)$; with $|\Phi|^2\le C^2_\Phi p^2\lambda_k^{-1/2}$ the standalone vertex is $O(t^{1/2}p^2\lambda^{-1/2})$ per cube times the tree factor. Every cube within $LM+3$ of $\partial\Omega^+_1$ lies in $B^k(\Lambda^{(k)}_7\setminus\Lambda^{(k)}_8)$ since $\dist(\Lambda_8,\Lambda^c_7)\ge r-1>LM+4$; for cubes in $B^k(\Lambda^{(k)}_8)$, $\dist(\square(x),\mathrm{Bad})\ge r-1-LM-3$. Count the cubes and use $p(t)^ne^{-\delta r(t)}\le t^\kappa$ for the exponential part.
\end{proof}

\begin{lemma}[the vacuum counterterm moves; ``even simpler'']\label{lem:a06d:D23}
$e_k\defeq[E_{k+1}-E_k]^{\mathrm{glob}}-[E_{k+1}-E_k]^{\mathrm{loc}}$ (localised with $\Omega^+_1$, $g^+$) is a $c$-number with $|e_k|\le Ct|T^{(k)}_1\setminus\Lambda^{(k)}_8|\le Ct(1+(17(r+M))^d)|\Lambda^{(k)c}_0|$; if $\Lambda^{(k)c}_0=\emptyset$ then $\mathrm{Bad}=\emptyset$ and $e_k=0$.
\end{lemma}

\begin{proof}
As in \cref{lem:a06d:D22} without external fields: Proposition~1 for vacuum families after cutting $S_{(k)}$ gives $O(t)$ per localisation (order $\ge2$). Root at a cube in $\mathrm{Bad}$, which lies in $T^{(k)}_1\setminus\Lambda^{(k)}_8$ up to the $LM+2$ layer, and sum. By \cite[(2.8)]{Ba2} every block of $\Lambda_0\setminus\Lambda_8$ is within $8(r+M)$ of $\Lambda^c_0$ (each layer is $r$ plus one large block of side $M$), and at most $(16(r+M)+1)^d\le(17(r+M))^d$ unit blocks lie within that distance of a given block of $\Lambda^c_0$.
\end{proof}

Both moves are bounded by $Ct^{1/2-\beta}(r+M)^d|\Lambda^{(k)c}_0|+C_\kappa t^\kappa|\Lambda^{(k)}_7|$, against which the large-field ledger of \cite[Sect.~3C]{Ba2} pays $(d+1)(1+\ln t^{-1})$ per unit of $|\Lambda^{(k)c}_0|$; that ledger is \cite[(3.42)]{Ba2}, part of the $L=5$ constants of \cref{thm:a06:BALA}.
